E_r(x,y) 
:= \sum_{i,s \ge 0} e_{i,(r^s)} x^i y^s 
= 1 + (1 + x) M_r(x,y),  
\]
where 
\[
M_r(x,y) 
:= \sum_{\substack{i \ge 0 \\ s \ge 1}} m_{i,(r^s)} x^i y^s.
\]
The product formula is a substantial merit of working with $e_{i,(r^s)}$, and it facilitates the extension of the case $s=1$ to $s>1$. 
This is done in Subsection~\ref{sec:nn}, where the result for $s=1$ is used to obtain the general case.
		
	
\subsection{Formulas for inner products}\label{sec:e}
	
In this subsection we obtain explicit formulas for the inner products 
\[
    e_{k,(r^s)} := \langle \psi^{(r^s)}, \chi^{(1^k) \oplus (n-k)} \rangle
    \quad (0 \le k \le n).
\]
Recall that, by Equation~\eqref{eq:alternating}, 
the sequences $\{m_k\}_{k=0}^{n-1}$ and $\{e_k\}_{k=0}^{n}$ 
determine each other, via the relations
\[
    e_k = m_k+m_{k-1}
    \text{ and }
    m_k = \sum_{i=0}^k (-1)^{k-i} e_i
    \qquad (0 \le k \le n),
\]
where $m_k := 0$ for $k = -1$ and $k = n$.
Nota bene, these multiplicities depend on $r$ and $s$
but this dependence is suppressed in the notation.

\begin{definition}\label{def:partition}
    For given non-negative integers $i$, $r$ and $s$, let
    \[
        P_{r,s}(i) 
        := \left\{ \gamma=(\gamma_1,\ldots,\gamma_s) :
        \sum_{\ell = 1}^{s} \gamma_\ell = i,\, r \ge \gamma_1 \ge \gamma_2 \ge \ldots \ge \gamma_s \ge 0 \right\}
    \]
    denote the set of all partitions of $i$ into at most $s$ parts, each of size at most $r$. 
    Denote the multiplicity of $j$ in $\gamma \in P_{r,s}(i)$ by 
    $k_j(\gamma) := |\{1 \leq \ell \leq s \mid \gamma_\ell=j\}|$.
\end{definition}
	
